\chapter{Chloride (Blood)}

\section*{What Is This Test?}
The chloride blood test measures the concentration of chloride ions (Cl\textsuperscript{-}) in the blood serum or plasma. Chloride is the major extracellular anion, accounting for approximately 70\% of the total anion content in extracellular fluid. Normal serum chloride ranges from 96--106 mEq/L (mmol/L). Chloride plays essential roles in maintaining osmotic pressure, fluid balance, and acid-base homeostasis. It is the primary counter-ion to sodium, so chloride levels generally parallel sodium levels in most clinical situations. Chloride is also critical for gastric acid (HCl) secretion by gastric parietal cells and for the chloride-bicarbonate exchange in red blood cells (the ``chloride shift'' or Hamburger phenomenon), which is essential for CO\textsubscript{2} transport. Renal chloride reabsorption occurs primarily in the proximal tubule, thick ascending limb of Henle (via Na-K-2Cl cotransporter, the target of loop diuretics), and in the distal tubule (via Na-Cl cotransporter, the target of thiazide diuretics). Chloride is routinely measured as part of the basic metabolic panel (BMP) and comprehensive metabolic panel (CMP), and is essential for calculating the anion gap [Anion Gap = Na -- (Cl + HCO\textsubscript{3})].

\section*{Alternative Names}
Cl; Serum Chloride; Plasma Chloride; Blood Chloride; Cl- Level; CSF Chloride (cerebrospinal fluid chloride measurement)

\section*{Principle}
Serum chloride is measured using ion-selective electrode (ISE) potentiometry with a chloride-sensitive membrane. The electrode contains a membrane impregnated with a quaternary ammonium salt or silver chloride that selectively binds chloride ions, generating an electrical potential proportional to chloride activity. On most automated analyzers, chloride is measured simultaneously with sodium and potassium on an integrated electrolyte panel using indirect ISE after sample dilution. Some instruments use coulometric titration with silver ions (generating silver chloride precipitate): Cl\textsuperscript{-} + Ag\textsuperscript{+} $\rightarrow$ AgCl, where the amount of silver ions required for complete precipitation is proportional to chloride concentration. Yet other methods use a mercuric thiocyanate colorimetric reaction: chloride displaces thiocyanate from mercuric thiocyanate; free thiocyanate reacts with ferric ions to form colored ferric thiocyanate measured spectrophotometrically at 480 nm. The ISE method is currently the most common in clinical laboratories. This test is not performed by hematology analyzers employing impedance or flow cytometry methods.

\section*{History}
Chloride was first recognized as a component of gastric juice by William Prout in 1824 (hydrochloric acid). The role of chloride in maintaining acid-base balance was established in the early 20th century through the work of Lawrence J. Henderson and Donald Van Slyke. The chloride shift mechanism in red blood cells was described by Hartog Jakob Hamburger in the 1890s. Colorimetric methods using mercuric thiocyanate were developed by Skeggs and Hochstrasser in the 1960s for the Technicon AutoAnalyzer. Ion-selective electrode potentiometry became the standard method in the 1970s-1980s. Photometric methods evolved from the late 19th century Volhard titration (silver nitrate with thiocyanate indicator) through Schales and Schales (1941) diphenylcarbazone-mercuric nitrate titration to current ISE technology.

\section*{Why This Test Is Ordered}
\begin{itemize}
  \item Part of the basic metabolic panel (BMP) and comprehensive metabolic panel (CMP) for routine health screening, preoperative assessment, and monitoring in hospitalized patients
  \item Evaluation and monitoring of acid-base disorders. Chloride is essential for calculating the anion gap [Na -- (Cl + HCO\textsubscript{3})], which is required to differentiate types of metabolic acidosis
  \item Assessment of volume status. Serum chloride trends often parallel sodium trends, and chloride deviations may signal water balance abnormalities
  \item Diagnosis of primary hyperchloremic (normal anion gap) metabolic acidosis, which is characterized by a low serum HCO\textsubscript{3} with normal anion gap (8--12 mEq/L) and elevated serum chloride
  \item Evaluation of persistent vomiting or nasogastric suction, leading to loss of gastric HCl and hypochloremic metabolic alkalosis (chloride-responsive metabolic alkalosis with urine chloride <10--20 mEq/L)
  \item Monitoring response to diuretic therapy. Loop diuretics (furosemide) and thiazide diuretics can cause significant chloride losses and subsequent hypochloremia
  \item Diagnosis of bromide toxicity. Bromide interferes with chloride measurement methods (especially ISE), causing apparent hyperchloremia that is not truly present---this can lead to a falsely elevated serum chloride and falsely low or negative anion gap
  \item Evaluation of adrenal disorders (Conn's syndrome, Cushing's syndrome, Addison's disease) and renal tubular acidosis
  \item Monitoring fluid resuscitation with normal saline (0.9\% NaCl), which contains 154 mEq/L of chloride (supraphysiologic) and can induce hyperchloremic metabolic acidosis with large volumes (>3--4 L)
  \item Workup of cystic fibrosis by sweat chloride testing (pilocarpine iontophoresis), although sweat chloride is a separate test measured on stimulated sweat, not serum
\end{itemize}

\section*{Is This a Primary or Secondary Test}
Chloride measurement is primarily a component of the electrolyte panel, making it both a primary screening test (as part of the BMP/CMP) and a secondary parameter, because isolated chloride abnormalities are rarely diagnostic alone and must be interpreted in the context of sodium, bicarbonate, anion gap, and acid-base parameters. Chloride is essential for the calculation of the anion gap, which is a cornerstone of metabolic acidosis evaluation.

\section*{How This Test Is Performed}
A healthcare professional draws blood from a vein using a small needle. After cleaning the skin with antiseptic and applying a brief tourniquet, blood is drawn into a serum separator tube (gold-top or red-top) or a lithium heparin tube (green-top). The sample volume needed is typically 3--5 mL. The sample is centrifuged to separate the serum or plasma from cells. Chloride is then measured on an automated chemistry analyzer using ISE potentiometry alongside other electrolytes. Results are available within 30 minutes to 1 hour. As with sodium and potassium, results from an EDTA (purple-top) tube are not appropriate for electrolyte measurement.

\section*{What Preparation Is Needed}
No special preparation or fasting is required for a serum chloride test. Because chloride generally parallels sodium, the same pre-analytical considerations apply: avoid excessive water intake immediately before the test, minimize prolonged tourniquet time during phlebotomy to avoid hemoconcentration, and do not draw from an arm receiving IV fluids (which would dilute or falsely elevate the result depending on the IV fluid composition). Inform the healthcare provider of all medications, particularly: diuretics (loop and thiazide), sodium bicarbonate therapy, corticosteroids, and laxatives. Bromide-containing medications (some older sedatives) can cause falsely elevated chloride readings. Prolonged tourniquet use can cause a slight increase due to hemoconcentration. The sample should be processed within 4 hours at room temperature.

\section*{Contraindications and Precautions}
No absolute contraindications for venipuncture. Important analytical interferences: (1) Bromide and iodide ions---both are sensed by chloride ISE with greater affinity than chloride itself, causing pseudohyperchloremia. This can result in a falsely low or even negative calculated anion gap. A negative anion gap should always trigger investigation for bromide toxicity. (2) Salicylate and thiocyanate are also highly sensed by some chloride electrodes and can cause falsely elevated chloride. (3) Lipemic samples---extreme hypertriglyceridemia can cause spurious chloride results with indirect ISE methods. (4) Drugs containing chloride salts (ammonium chloride, arginine hydrochloride, potassium chloride) can cause true hyperchloremia. (5) Hemolysis does not significantly affect chloride measurement. (6) Volume depletion/dehydration causes a proportional increase in chloride along with sodium. (7) Large-volume normal saline resuscitation (0.9\% saline) contains supraphysiologic chloride (154 mEq/L vs. normal plasma ~103 mEq/L) and can cause hyperchloremic metabolic acidosis. (8) Recent consumption of a high-salt meal has minimal effect on serum chloride in individuals with normal renal function.

\section*{Risks}
Minimal risk associated with venipuncture. Mild pain or bruising at the site is common and typically resolves within 1--2 days. Rare complications include hematoma, infection at the puncture site, inadvertent arterial puncture, nerve injury, and vasovagal syncope (fainting). No risk of significant blood loss (3--5 mL drawn). No additional risk specific to chloride measurement beyond standard venipuncture.

\section*{What Happens After the Test}
A cotton ball and bandage are applied over the puncture site with gentle pressure for 1--2 minutes to prevent hematoma. Results are usually available within 30--60 minutes. The chloride value alone is rarely interpreted in isolation; it is always evaluated alongside sodium, potassium, bicarbonate, and the calculated anion gap. Specific clinical scenarios: (1) If chloride is low with low bicarbonate (metabolic acidosis with normal anion gap), further investigation for hyperchloremic metabolic acidosis is warranted. (2) If chloride is very low (<90 mEq/L) with metabolic alkalosis and low urine chloride (<10 mEq/L), this suggests chloride-responsive metabolic alkalosis (e.g., vomiting, NG suction, diuretic use) that will respond to volume and chloride repletion with normal saline. (3) If chloride is elevated with a low or negative anion gap, pseudohyperchloremia from bromide should be considered. (4) If chloride tracks proportionally with sodium changes, the primary disorder is likely a water balance problem rather than an isolated chloride disorder.

\section*{Understanding Results}

\subsection*{Below Normal}
\begin{itemize}
  \item \textbf{Hypochloremia (Cl <96 mEq/L):} Often part of broader electrolyte abnormalities and rarely isolated.
  \item \textbf{Gastrointestinal losses (most common):} Persistent vomiting; nasogastric (NG) suction---loss of HCl-rich gastric fluid causes metabolic alkalosis with hypochloremia and hypokalemia. With NG suction, each liter of gastric fluid contains approximately 120--160 mEq of chloride. Severe diarrhea can also lower chloride.
  \item \textbf{Renal losses:} Loop diuretic therapy (furosemide, bumetanide, torsemide) blocks the Na-K-2Cl cotransporter in thick ascending limb of Henle; thiazide diuretics block distal Na-Cl cotransporter. Both cause urinary chloride loss. Chronic metabolic alkalosis from diuretic use. Salt-wasting nephropathy. Post-obstructive diuresis. Recovery phase of acute tubular necrosis. Bartter syndrome (genetic defect in thick ascending limb transporters) and Gitelman syndrome (defect in distal Na-Cl cotransporter).
  \item \textbf{Metabolic alkalosis with low urine chloride (<10--20 mEq/L):} Chloride-responsive metabolic alkalosis indicating volume depletion and need for chloride repletion. Causes: vomiting, NG suction, post-hypercapnia alkalosis, diuretic use (remote). This responds to normal saline (NaCl) administration.
  \item \textbf{Metabolic alkalosis with high urine chloride (>20 mEq/L):} Chloride-resistant metabolic alkalosis; indicates mineralocorticoid excess or ongoing diuretic therapy. Causes: primary hyperaldosteronism, Cushing's syndrome, Bartter/Gitelman syndromes, licorice ingestion, severe hypokalemia. Does not respond to saline.
  \item \textbf{Chronic compensatory respiratory acidosis:} Renal compensation for chronic CO\textsubscript{2} retention increases bicarbonate reabsorption and reciprocally lowers chloride.
  \item \textbf{SIADH:} Water retention causes dilution of both sodium and chloride proportionally.
  \item \textbf{Adrenal insufficiency (Addison's disease):} Hyponatremia and hypochloremia.
  \item \textbf{Congestive heart failure or cirrhosis:} Dilutional hyponatremia/hypochloremia from water retention, often with restrictive fluid intake.
  \item \textbf{Extensive burns:} Loss of chloride-rich fluid through damaged skin.
\end{itemize}

\subsection*{Above Normal}
\begin{itemize}
  \item \textbf{Hyperchloremia (Cl >106 mEq/L):} Most commonly iatrogenic or associated with hyperchloremic metabolic acidosis (normal anion gap).
  \item \textbf{Hyperchloremic (normal anion gap) metabolic acidosis:} Characterized by low serum bicarbonate with normal anion gap (8--12 mEq/L) and elevated chloride. Causes: (a) GI bicarbonate loss---severe diarrhea (loss of bicarbonate-rich pancreatic and intestinal secretions), pancreatic fistula, ileostomy drainage, cholestyramine (binds bile acids that normally stimulate bicarbonate secretion), ureteral diversion (ureterosigmoidostomy---urinary chloride exchanged for bicarbonate in colon); (b) Renal tubular acidosis (RTA)---Type I (distal RTA: defect in distal H+ secretion, seen in Sj\"{o}gren's syndrome, amphotericin B, autoimmune), Type II (proximal RTA: defect in proximal HCO\textsubscript{3} reabsorption, seen in Fanconi syndrome, multiple myeloma, carbonic anhydrase inhibitors like acetazolamide), Type IV (hyporeninemic hypoaldosteronism: most common type, seen in diabetic nephropathy, ACE inhibitor/ARB use, trimethoprim); (c) Administration of chloride-rich solutions---large-volume 0.9\% normal saline resuscitation (154 mEq/L Cl, pH ~5.5---causes dilutional acidosis and hyperchloremia), total parenteral nutrition (TPN) with excessive chloride from amino acid solutions, ammonium chloride administration; (d) Early recovery from diabetic ketoacidosis---ketone anions are metabolized to bicarbonate; if chloride replacement was excessive during treatment, hyperchloremia persists.
  \item \textbf{Dehydration/hypovolemia:} Proportional increase in chloride along with sodium from water loss (fever, sweating, inadequate water intake, diabetes insipidus). Serum chloride and sodium rise proportionally.
  \item \textbf{Primary hyperparathyroidism:} Can cause hyperchloremic metabolic acidosis through the effect of PTH on renal bicarbonate handling (PTH inhibits proximal tubular bicarbonate reabsorption).
  \item \textbf{Acetazolamide therapy:} Carbonic anhydrase inhibitor that blocks proximal tubular bicarbonate reabsorption; causes bicarbonate loss with compensatory chloride retention, resulting in hyperchloremic metabolic acidosis.
  \item \textbf{Bromide toxicity/pseudohyperchloremia:} Bromide ion has greater affinity for the chloride ISE (selectivity coefficient of ~2.5:1 relative to chloride), producing artificially high chloride measurements. This creates a low or negative anion gap. Historical bromism from sedatives (e.g., Bromo-Seltzer, now rare) but still relevant in patients using some herbal or imported medications.
  \item \textbf{Renal failure:} In advanced CKD, reduced bicarbonate production leaves chloride as the dominant anion, with a proportional increase in serum chloride.
\end{itemize}

\section*{Normal Results}
\begin{itemize}
  \item \textbf{Adults:} 96--106 mEq/L (mmol/L)
  \item \textbf{Children (1--16 years):} 98--107 mEq/L
  \item \textbf{Infants (1--12 months):} 96--108 mEq/L
  \item \textbf{Newborns (0--7 days):} 96--110 mEq/L
  \item \textbf{CSF chloride:} 118--132 mEq/L (higher than serum)
  \item \textbf{Sweat chloride:} <30 mmol/L (normal); 30--59 mmol/L (borderline); $\geq$60 mmol/L (diagnostic for cystic fibrosis)
\end{itemize}

\section*{Follow-up Tests}
\begin{itemize}
  \item \textbf{Basic metabolic panel or comprehensive metabolic panel:} Complete electrolyte assessment including sodium, potassium, bicarbonate, BUN, creatinine, glucose, anion gap to contextualize chloride abnormalities.
  \item \textbf{Anion gap calculation:} AG = Na -- (Cl + HCO\textsubscript{3}); normal 8--12 mEq/L. Required to classify metabolic acidosis as high-AG or normal-AG (hyperchloremic).
  \item \textbf{Arterial blood gas (ABG):} To assess acid-base status, pH, pCO\textsubscript{2}, and calculated bicarbonate. Essential when chloride is abnormal to determine if metabolic acidosis or alkalosis is present.
  \item \textbf{Urine chloride:} Measurement of urine chloride concentration (spot or 24-hour). Urine chloride <10--20 mEq/L in metabolic alkalosis suggests chloride-responsive alkalosis (volume depletion, vomiting, NG suction). Urine chloride >20 mEq/L suggests chloride-resistant alkalosis (mineralocorticoid excess, Bartter/Gitelman, active diuretic therapy).
  \item \textbf{Serum osmolality and osmolal gap:} If toxic alcohol ingestion is suspected in setting of high-AG metabolic acidosis (which may coexist with chloride abnormalities).
  \item \textbf{Plasma renin activity and aldosterone:} If mineralocorticoid excess is suspected as cause of hypochloremic metabolic alkalosis.
  \item \textbf{24-hour urine electrolytes:} Comprehensive assessment if Bartter or Gitelman syndrome is considered.
  \item \textbf{Sweat chloride test (pilocarpine iontophoresis):} For suspected cystic fibrosis if family history or clinical features suggest the diagnosis.
  \item \textbf{Serum bromide level:} If pseudohyperchloremia with negative or very low anion gap suggests bromide toxicity.
\end{itemize}

\section*{References}
\begin{enumerate}
  \item Burtis CA, Ashwood ER, Bruns DE. Tietz Textbook of Clinical Chemistry and Molecular Diagnostics. 7th ed. Elsevier; 2023.
  \item Wallach J. Interpretation of Diagnostic Tests. 10th ed. Wolters Kluwer; 2022.
  \item Tierney LM, Saint S, Drazen JM. CURRENT Medical Diagnosis and Treatment. McGraw-Hill; 2023.
  \item Fischbach FT, Dunning MB. A Manual of Laboratory and Diagnostic Tests. 10th ed. Wolters Kluwer; 2023.
  \item Gennari FJ, Weise WJ. Acid-base disturbances in gastrointestinal disease. Clin J Am Soc Nephrol. 2008;3(6):1861--1868.
  \item Seifter JL. Integration of acid-base and electrolyte disorders. N Engl J Med. 2014;371(19):1821--1831.
  \item Galla JH. Metabolic alkalosis. J Am Soc Nephrol. 2000;11(2):369--375.
  \item Kraut JA, Madias NE. Serum anion gap: its uses and limitations in clinical medicine. Clin J Am Soc Nephrol. 2007;2(1):162--174.
\end{enumerate}

\clearpage
